\section{The Year Still Given}

At the year's end the Church's expectation of Christ's return brings the whole life back into view. The months have contained duties fulfilled and neglected, graces welcomed and resisted, losses endured, and opportunities that will not return. Some matters call for thanksgiving; some for confession; some for patient amendment. Meditation on judgment becomes fruitful when it allows these realities to be named before God without either evasion or despair.

We need not devise an extraordinary life before we can become ready. The neglected person may already be known to us. The sin that needs confession may already have been recognized. The restitution delayed, the prayer habit abandoned, and the word of forgiveness withheld may have occupied conscience for a long time. The Lord's coming gives these familiar things their true urgency. We can begin with the next act of fidelity rather than waiting for a moment in which conversion will cost less.

The year still given allows us to act on what examination has made clear. If confession has been delayed, we can arrange to go. If an injury needs repair, we can seek counsel about making it. A duty of care can be resumed before the familiar excuse returns. We undertake these acts with hope in Christ, who calls us to repentance and gives the grace to answer him.

